\section{Curvature reads under one actual Gaussian reader}
\label{sec:actual-reader-guarded-curvature}

This section records the nonlinear step between a local Gaussian reader and
the group-valued field that it is meant to read. All comparisons use one
coarse field, one gauge based at one site, one parent forecast, and one joint
Gaussian sample. A patch restricts which addresses are tested; it does not
change the forecast, covariance, or gauge from address to address. The
local Gaussian hypotheses used here are those of
Section~\ref{sec:actual-reader-local-gaussian-prices}. That section's
same-pose mean and covariance estimates remain inputs, as do the paid
regularity and window hypotheses specified below.

Let $x$ be the large coarse parameter, $L$ the fixed blocking side, and
\[
 y=x+c_L,\qquad R=x^{-3/8},\qquad
 t=\operatorname{phase}(y),\qquad
 g=\operatorname{coupling}(\beta),\quad
 g'=\operatorname{coupling}(\beta').
\]
Here $c_L$ is the fixed scale shift, $x/2\leq\beta\leq3x$,
$x/2\leq\beta'$, and $|\beta'-\beta|\leq\gamma$. The change from $\beta$
to $\beta'$ is paid by the relative-coupling allowance $x^{-1/8}$. Write
$\tau=\operatorname{phase}(x)$ for the coarse phase and
$t=\operatorname{phase}(y)$ for the fine phase; the scale receipt is
$\tau\leq2t$. The scalar functions have the
definitions
\[
 R=\operatorname{linkRoom}(x),\qquad
 \operatorname{walls}(y)=\lceil(\log y)^2\rceil^2,\quad
 \operatorname{phase}(y)=y^{-1/2}
       \exp((\log\log y)^{3/2}).
\]
Let $C_{\rm core}=\operatorname{priceCore}(G)$ be the fixed price for the
coarse response chart, $C_{\rm high}=\operatorname{highPrice}(G)$ the fixed
price for a fine link in the exponential chart, and let
$A_*$ be the fixed price for the two-link curvature remainder, chosen at least
as large as $\operatorname{arcPrice}(G)$. Let $n_G$ be the real color dimension and set
\[
 e=(A_*+n_G)R^2.
\]
These are distinct prices and each estimate below uses its stated price;
we do not identify the chart, link, and curvature constants. Choose them
before the large-$x$ threshold and require the eventual heat and transport
inequalities for these fixed values.
Let $\mathcal C_o(V)_q$ denote the nonlinear plaquette-curvature coordinate
of the coarse field in the fixed gauge based at $o$, indexed by coarse
plaquettes and color coordinates. The parent forecast is the fixed real row
\[
 \mathcal P_{o,V,g'}(u,a)=\frac1{g'}
   \sum_q K_{u,a;q}\,\mathcal C_o(V)_q.
\]
The finite linear response operator selected by the patch and forest will be
denoted $D_{A,F,x}$; it acts on the one joint Gaussian vector and retains
the source transport and gauge data.
When an incidence index $j$ is used instead of $(u,a)$, it denotes that
same plaquette/color row. The same $V,o,g'$ and kernel $K$ are used for every tested row. In
particular, $\mathcal P$ is not recomputed from a patch-dependent field.

\subsection{The scalar scale and ticket estimates}

Fix a patch of side $L$, a finite transition tube, its paid row and chart
bounds, a fixed width $w$, an action budget $s>0$, a positive slope allowance
$\gamma$, and a nonnegative intensity $z$. There is a threshold $x_0$ such
that for every $x\geq x_0$ and every $\beta,\beta'$ in the bands above,
the following scalar estimates hold. Set
\[
 M_0=\frac{\operatorname{psiBudget}(z)}{L^2},
 \qquad P_0=6M_0=|\mathrm{TypeFace}_2|M_0=\operatorname{ticket}(z,L),
\]
Here $\mathrm{TypeFace}_2$ is the set of two-element subsets of the four
coordinate axes, so its cardinality is six. The separate factor $2$ in the
room budget below comes from $\tau\leq2t$. By taking the fixed $L$ large
enough, source sizing gives
$2n_GP_0\leq1/4096$. For this $L$, the eventual source estimates give
\begin{align}
 &\operatorname{clearance}(x)\leq R\leq\operatorname{pinMargin}
   \quad\text{and}\quad4C_0\operatorname{clearance}(x)\leq R,
   &x^{-7/8}&\leq R^2, \label{eq:readerB-clearance}\\
 &\operatorname{docCost}(x)\leq C_0 R,
   \quad \operatorname{cheque}(s,M_0,x)\leq R,
   &t\leq 2x^{-7/16}, \label{eq:readerB-chart}\\
 &g^{-1}\leq2\sqrt{x},\qquad
 &\tau&\leq2t,\label{eq:readerB-coupling}\\
 &C_{\rm core}R^2\,\operatorname{stride}_4(x)&=o(x^{-1/8}),
   &e/g&=O(x^{-1/4}).\label{eq:readerB-power}
\end{align}
The second clearance estimate is the scale-scope bound applied with
multiplier $4C_0$; it is needed to turn the chart estimate into a
quarter-room bound. The scalar ticket bound gives
\[
 \operatorname{retract}(x)
   =\operatorname{cheque}(s,M_0,x)C_{\rm core}R
       +P_0C_{\rm core}R^2\leq C_{\rm core}(1+P_0)R^2.
\]
The relative-coupling part is bounded by
$P_0\tau\,\operatorname{coupler}(x)$. These are the two terms paid by the one-step budget
\begin{equation}\label{eq:readerB-ticket-bill}
 \frac{\operatorname{retract}(x)+P_0\tau\,\operatorname{coupler}(x)}{g}
       \leq x^{-1/16}.
\end{equation}
Here the relative-coupling receipt is
$|g'^{-1}g-1|\leq\operatorname{coupler}(x)$.
The remaining large-$x$ comparisons used below are
\begin{align}
 4C_{\rm high}R&\leq\delta_G,\label{eq:readerB-small-chart}\\
 6c_G(C_{\rm high}R)^2+4e&\leq t/64,\label{eq:readerB-four-budget}\\
 \operatorname{nautical}(G,C_{\rm high}R,g,w\operatorname{walls}(y))
       +4e/g&\leq x^{-1/8}.\label{eq:readerB-nautical}
\end{align}
Here $c_G$ is the fixed color-scale constant and $\delta_G$ is the
near-identity radius in the quadratic curvature lower bound.
For the transport term, the source formula is explicit. With
$M=C_{\rm high}R$, $B=w\operatorname{walls}(y)$, and $u=4BM$,
\[
 \operatorname{nautical}(G,M,g,B)
 =\frac{c_G\{u(4M)+(u+4M)u+6M^2\}}{g}.
\]
For fixed $w$ and $y=x+c_L$, the wall definition gives both
$B=O(\operatorname{stride}_4(x)^{1/2})$ and
$B^2=O(\operatorname{stride}_4(x))$. Since the displayed polynomial in $B$
is quadratic and $M=C_{\rm high}R$, the right side is at most
$C\operatorname{stride}_4(x)R^2/g=O(\operatorname{polylog}(x)x^{-1/4})$.
These estimates follow with the stated powers: $R=x^{-3/8}$,
$t\leq2x^{-7/16}$, $g^{-1}=O(\sqrt{x})$, and
$\operatorname{walls}(y)\leq\operatorname{stride}_4(x)$; the local chart
and ticket prices fix the constants before $x$ varies. In particular
$R^2= x^{-3/4}=o(t)$ because
$t\geq y^{-1/2}$, and $e/g=O(x^{-1/4})$.

The quadratic lower bound used here is the elementary local curvature
estimate: whenever the group chord of $U$ is at most $\delta_G$,
\[
 \frac34\,\operatorname{chord}(U)^2
       \leq\|\kappa_G(U)\|^2.
\]
Its constants depend only on the fixed compact gauge group. The paid
chart, action, and ticket inputs are exactly the scalar hypotheses consumed
by the local reader argument; no uniformity in $L$, patch, or tube is
claimed here.

\begin{lemma}[Eventual local room]
\label{lem:readerB-enlightened}
For fixed $G,L,w,z,s,\gamma$ as above and a transition tube with the stated
paid chart bounds, all sufficiently large $x$ and all
$x/2\leq\beta\leq3x$ have the following properties:
\begin{enumerate}
\item the chart clearance is at most $R$, $R$ lies below the fixed chart
   margin, and $x^{-7/8}\leq R^2$;
\item every funded fluctuation of norm at most the clearance has deterministic
   chart size plus $n_Gx^{-7/8}$ and the Gaussian reserve
   $\kappa_0t$ at most $R$, where $\kappa_0=1/4096$;
\item the near-identity link price is at most $C_{\rm high}R$, and the four-face
   curvature budget is at most $t/64$;
\item the sum of the four-face gauge-reading error, the coupling error and
   the residual curvature price is at most $x^{-1/32}$.
\end{enumerate}
The assertions hold uniformly in the period for the fixed patch and local
data.
\end{lemma}
\begin{proof}
The selected row and chart receipts bound the deterministic logarithmic
response by $C_0$ times the available chart clearance. The source
$4C_0\operatorname{clearance}(x)\leq R$ makes this at most $R/4$.
Since
$x^{-7/8}/R=x^{-1/2}$ and
$t/R\leq2x^{-1/16}$, the color-count and reserve terms are each at most
$R/4$ after increasing $x_0$. This proves the room estimate. The link price
is the fixed chart constant times $R$; the four-face budget is
$O(R^2)=O(x^{-3/4})$, while $t\geq y^{-1/2}\geq c x^{-1/2}$ for large
$x$, proving the curvature budget. Finally, the definition of the
four-face gauge error is
\[
 \operatorname{nautical}(G,C_{\rm high}R,g,w\operatorname{walls}(y))
          +4e/g.
\]
The scalar scale estimate bounds this by $x^{-1/8}$. Adding the paid row
read error $x^{-1/16}$ and the remaining coupling allowance gives at most
$x^{-1/32}$ eventually, since each stronger negative power divided by
$x^{-1/32}$ tends to zero. The fixed costs were chosen before the period,
so the threshold is period-independent.
\end{proof}

\subsection{From the coarse forecast to a fine curvature row}

Let the coarse side be $c$ and the fine side be $f$, with a fixed blocking
map. Let $V$ be the coarse field and $o$ a single base site. The common pose
assumption is that one based gauge transforms $V$ on the whole room about
$o$ into the exponential of one real direction $v$, with
$\|v\|\leq\operatorname{clearance}(x)/2$. The field is healthy on that
room, and the chosen stores have no support there. The same gauge is used
for every row and plaquette below.

Let $F$ be the selected forest and $Z$ the actual joint Gaussian sample on
its active link coordinates. Write $\ell$ for the low, linear response to
the coarse direction, and write $\mathrm{doc}$ for the paid deterministic
local response. Write $\mathcal G_{A,F,x}(Z)=D_{A,F,x}Z$ for the transported Gaussian
response. The fine field on an interior bond $j$ has the exact form
\begin{equation}\label{eq:readerB-fine-link}
 U_j=\exp(w_j+g z_j)\,\operatorname{replica}(V^o)_j,
 \qquad V^o=\operatorname{spine}(V,o),
\end{equation}
where $w_j$ is the local portrait read and $z_j$ is the corresponding
Gaussian coordinate combination. Write $d_j$ for the deterministic
response, $b_j$ for the logarithm of the replica link, and $r_j=gz_j$.
The source identities give $\ell=d+b$ and
$\|w_j-d_j\|\leq n_Gx^{-7/8}$. The actual link is $\exp(w_j+r_j)\exp(b_j)$. Write $\mathsf M$ for the
fixed unitary matrix representation, viewed as a real vector space, and
$\Pi_{\mathfrak g}$ for the normalized real-linear color projection onto
the Lie algebra, with operator norm at most $1/c_G$. The curvature
coordinate is
\[
 \kappa_G(U)=\Pi_{\mathfrak g}\!\left(
       \frac{\mathsf M(U)-\mathsf M(U^{-1})}{2}\right).
\]
For a unitary group element this also equals
$\Pi_{\mathfrak g}(\mathsf M(U)-I)$. The latter identity gives the exact
two-factor relation
\[
 \kappa_G(UV)-\kappa_G(U)-\kappa_G(V)
   =\Pi_{\mathfrak g}((\mathsf M(U)-I)(\mathsf M(V)-I)).
\]
The Frobenius product inequality and
the bound
\[\|\mathsf M(U)-I\|_F\leq c_G\operatorname{chord}(U)\]
bound this defect by $c_G\operatorname{chord}(U)\operatorname{chord}(V)$. Also use the radial
chart estimate $\|\kappa_G(\exp a)-a\|\leq A_{\rm arc}\|a\|^2$
for $\|a\|$ within the fixed chart, with $A_{\rm arc}=\operatorname{arcPrice}(G)$.
Since $\operatorname{chord}(\exp a)\leq\|a\|$, for $U=\exp a$ and
$V=\exp b$ we obtain
\begin{align*}
 \|\kappa_G(UV)-(a+b)\|
 &\leq c_G\|a\|\|b\|+A_{\rm arc}(\|a\|^2+\|b\|^2)\\
 &\leq (c_G/4+A_{\rm arc})(\|a\|+\|b\|)^2.
\end{align*}
Choose $A_*\geq c_G+4A_{\rm arc}$. This is the stated curvature-coordinate
bound, not a logarithmic estimate.
Together with $x^{-7/8}\leq R^2$ it absorbs the portrait error into the
fixed quadratic price. The guarded event is
\[
 \|g z_j\|\leq\kappa_0t
\]
for every bond in the enlarged room. The paid room estimate is
\[
 \|w_j+r_j\|\leq \|d_j\|+n_Gx^{-7/8}+\kappa_0t\leq R,
\]
and the blueprint estimate gives $\|b_j\|\leq R$. The source's two-link
chart estimate, proved by Taylor expansion of the curvature coordinate at
the identity, now gives
\begin{align}
 \operatorname{chord}(U_j)&\leq C_{\rm high}R,\label{eq:readerB-bondchord}\\
 \|\kappa_G(U_j)-(\ell_j+r_j)\|&\leq e.
 \label{eq:readerB-bondcurv}
\end{align}
Indeed, $w_j+r_j+b_j=(d_j+b_j)+r_j+(w_j-d_j)=\ell_j+r_j+(w_j-d_j)$.
The portrait term is at most $n_Gx^{-7/8}$; since both Taylor arguments
have norm at most $R$, the two-link curvature remainder is at most $A_*R^2$. Since $x^{-7/8}\leq R^2$, their sum is
bounded by $e=(A_*+n_G)R^2$. Here $\ell_j+gz_j$ is the value of the link direction at the single bond
$j$; the face boundary sum is taken only in the next subsection.

\begin{lemma}[Read error against the common parent forecast]
\label{lem:readerB-pointing}
Assume the paid row and transport estimates for the selected transition
tube: on every readable address in the room, the fine response row differs
from its shorter parent row by total absolute mass at most $\Delta$, the
parent row has mass at most $P_0$, and the deterministic coarse curvature
Taylor bounds stated below hold. Let
$\Delta=\operatorname{cheque}(s,M_0,x)$, the source fee at base price
$M_0=\operatorname{psiBudget}(z)/L^2$. Then for every interior fine
plaquette $u$ and color coordinate $a$, with
$f=\operatorname{faceCurl}(u,\ell)$,
\begin{equation}\label{eq:readerB-pointing-conclusion}
 \|f\|\leq t/2048,\qquad
 \left|g^{-1}f_a-\mathcal P_{o,V,g'}(u,a)\right|
      \leq x^{-1/16}.
\end{equation}
\end{lemma}
\begin{proof}
Index coarse plaquette/color coordinates by $q$. Let $m_q$ be the
linearized coarse plaquette curvature in direction $v$, and let
$(u_0)_q=\mathcal C_o(V)_q$ be the nonlinear curvature in the same based
gauge. The exact response identities are
\[
 f_a=\sum_q P_{u,a;q}m_q,
 \qquad Q_{u,a}=\sum_qK_{u,a;q}(u_0)_q,
 \qquad \mathcal P_{o,V,g'}(u,a)=Q_{u,a}/g'.
\]
The deterministic chart estimates give the global linear-row bound
$|m_q|\leq C_{\rm core}R$ for every tested coarse plaquette/color
coordinate. On the support of the parent row $K_{u,a;\cdot}$ they also give
$|m_q-(u_0)_q|\leq C_{\rm core}R^2$.
The purchased row bounds are
\[
 \sum_q|P_{u,a;q}-K_{u,a;q}|\leq\Delta,
 \qquad \sum_q|K_{u,a;q}|\leq P_0.
\]
Subtract the response sums and add/subtract $\sum_qK_{u,a;q}m_q$. Then
\[
 |f_a-Q_{u,a}|
 \leq C_{\rm core}R\sum_q|P_{u,a;q}-K_{u,a;q}|
       +C_{\rm core}R^2\sum_q|K_{u,a;q}|
 \leq \Delta C_{\rm core}R+P_0C_{\rm core}R^2
 =\operatorname{retract}(x).
\]
For the normalized comparison, $|Q_{u,a}|\leq P_0\tau$ since the
nonlinear parent curvature coordinates have modulus at most $\tau$. Thus
\[
 \left|\frac{f_a}{g}-\frac{Q_{u,a}}{g'}\right|
 \leq\frac{|f_a-Q_{u,a}|+|Q_{u,a}|\,|g'^{-1}g-1|}{g}
 \leq\frac{\operatorname{retract}(x)+P_0\tau\operatorname{coupler}(x)}g
 \leq x^{-1/16},
\]
by the relative-coupling receipt and \eqref{eq:readerB-ticket-bill}. For
the vector bound, the same $|Q_{u,a}|\leq P_0\tau$ and
$|f_a-Q_{u,a}|\leq\operatorname{retract}(x)$ hold for every color
coordinate. The purchased torque inequality is
\[
 n_G(P_0\tau+\operatorname{retract}(x))\leq t/2048,
\]
which gives $\|f\|\leq t/2048$ in the source color norm. All rows have
finite support in the selected window, and the bounds are period-uniform.
\end{proof}

\subsection{The four-bond curvature and the common gauge read}

For a fine plaquette $u$, let its four oriented bonds be
$j_1,j_2,j_3,j_4$, with signs $\epsilon_i$ determined by orientation, and
let $\Xi_u=U_{j_1}U_{j_2}U_{j_3}^{-1}U_{j_4}^{-1}$. The face curl
$\operatorname{faceCurl}_u(A)$ is the corresponding signed sum of the four
link-direction coordinates $A_{j_i}$.
Apply the two-factor relation above successively to the four oriented
factors $U_{j_1},U_{j_2},U_{j_3}^{-1},U_{j_4}^{-1}$. Each has chord at most
$M=C_{\rm high}R$, and inversion preserves chord while negating curvature.
At the second factor the defect is at most $c_GM^2$. At the third, the
preceding product has chord at most $2M$, so the new defect is at most
$2c_GM^2$. At the fourth, its chord is at most $3M$, so the new defect is
at most $3c_GM^2$. Summing gives exactly
$(1+2+3)c_GM^2=6c_GM^2$. Replacing the four
individual link-curvature coordinates by their linear directions using
\eqref{eq:readerB-bondcurv} costs at most $4e$. Consequently
\begin{equation}\label{eq:readerB-four-product}
 \|\kappa_G(\Xi_u)-\operatorname{faceCurl}_u(\ell+gZ)\|
       \leq 6c_G(C_{\rm high}R)^2+4e\leq t/64,
\end{equation}
by \eqref{eq:readerB-four-budget}. Linearity of the face curl, the bound
from Lemma~\ref{lem:readerB-pointing}, and the four guarded link increments
give
\[
 \|\operatorname{faceCurl}_u(\ell+gZ)\|
       \leq t/2048+4\kappa_0t.
\]
Consequently the product curvature coordinate has norm at most
$t/2048+4\kappa_0t+t/64<t/32$.
Since $\kappa_0=1/4096$,
\[
 \frac1{2048}+\frac4{4096}+\frac1{64}<\frac1{32}.
\]
The plaquette chord is at most four times the link chord, so
$4C_{\rm high}R\leq\delta_G$ keeps it inside the local curvature chart.
The local quadratic curvature lower bound therefore converts the
curvature norm bound into the required group-chord bound, with the source
normalization giving
\begin{equation}\label{eq:readerB-healthy-face}
 \operatorname{chord}(\Xi_u)\leq t/8.
\end{equation}
Thus if the health flag threshold is the phase $t$, the tested flag is
strictly below threshold. Applying this to each plaquette in the tested
patch proves the healthy-mask conclusion. It is a local conclusion on the
guarded sample; it does not assert that the field is healthy outside the
specified room.

The final read uses the same patch gauge as the forecast. If $U'$ is any
fine field with a gauge transform equal to $U$ on the patch, the exact
four-corner identity gives
\begin{multline*}
 \left|\operatorname{rigid}
       \bigl(\operatorname{coloring}(f'(r))
              (\operatorname{kiosk}_g(r,U'))\bigr)(j)
      -\bigl(\mathcal P_{o,V,g'}(j)
       +\mathcal G_{A,F,x}(Z)(j)\bigr)\right|\\
 \leq \operatorname{nautical}(G,C_{\rm high}R,g,w\operatorname{walls}(y))
       +4e/g+x^{-1/16}.
\end{multline*}
The first two terms are the four-link gauge transport error: the corner
product is compared with the four curvature rows and each row costs at most
$e/g$. The last term is the parent-forecast discrepancy from
\eqref{eq:readerB-pointing-conclusion}. By
\eqref{eq:readerB-nautical}, this is at most
$x^{-1/8}+x^{-1/16}\leq x^{-1/32}$ for large $x$. The gauge can vary in
this final comparison, but the forecast, sample, and original pose do not.

\begin{proposition}[Actual-reader curvature and read estimate]
\label{prop:readerB-curvature-read}
Fix a selected transition tube and one posed coarse field with based gauge
at $o$. Suppose the actual Gaussian reader estimates of
Section~\ref{sec:actual-reader-local-gaussian-prices}, the paid scalar
estimates above, and the local window/separation conditions hold. For each
room-view sample whose Gaussian coordinates are below the guard threshold,
the resulting fine field is healthy on the selected patch. For every gauge-equivalent fine field and every incidence row $j=(u,a)$
whose two-step neighborhood lies inside the patch, its rigid color read differs from the same parent
forecast plus the transported Gaussian curvature by at most
$x^{-1/32}$. All conclusions are uniform in the period once the displayed
fixed-data threshold is reached.
\end{proposition}
\begin{proof}
Use \eqref{eq:readerB-fine-link} for all bonds in the enlarged room, then
apply Lemma~\ref{lem:readerB-pointing} to each tested row. The four-bond calculation gives \eqref{eq:readerB-healthy-face}; the four-corner identity
and \eqref{eq:readerB-nautical} give the read estimate. All statements use
the same base site and parent forecast, and their constants do not depend
on the number of sites outside the tested patch.
\end{proof}

\subsection{The lost-mask estimate for a correlated Gaussian}

Let $\mathcal I_y$ be the actual finite set of real link-color coordinates
read by the guard on the enlarged room. Its cardinality satisfies
\begin{equation}\label{eq:readerB-room-dimension}
 m_y:=|\mathcal I_y|\leq C_w(\log y)^{16}
\end{equation}
for a constant $C_w$ depending on the fixed width and color dimension, but
not on the period. Indeed the room radius is $w\operatorname{walls}(y)+8$;
$\operatorname{walls}(y)\leq4(\log y)^4$ for large $y$, and the number of
link-color coordinates in a four-dimensional box is at most a fixed
multiple of its side to the fourth power.

If $m_y=0$, the product veil is $1$ and the loss is zero. Otherwise assume
$m_y\geq1$; the gauge algebra is nontrivial, so $n_G\geq1$. Group the
coordinates by bond, writing $X_b\in\mathfrak g$ for the actual Lie-valued
fluctuation read on bond $b$. Let $X=(X_i)_{i\in\mathcal I_y}$
have the actual centered joint Gaussian law. Its covariance matrix is
positive semidefinite and each entry has
absolute value at most $U$, where $U\geq1$ is the common palette bound from
the actual reader estimates. For $v>0$, each marginal obeys the Chernoff
bound
\[
 \mathbb P(|X_i|\geq v)\leq2\exp(-v^2/(2U));
\]
this follows by optimizing
$\mathbb E e^{sX_i}=e^{s^2\operatorname{Var}(X_i)/2}$ at $s=v/U$ and
applying the same estimate to $-X_i$. A union bound, which does not assume
independence, gives
\[
 \mathbb P\{\max_i|X_i|\geq v\}
       \leq2m_y\exp(-v^2/(2U)).
\]
The source veil is a product over bonds, and its away event is
$\max_b\|X_b\|\geq v$. If this occurs, one of the $n_G$ real coordinates
of some $X_b$ has modulus at least $v/\sqrt{n_G}$. A union bound therefore
gives
\[
 \mathbb P\{\max_b\|X_b\|\geq v\}
       \leq2m_y\exp(-v^2/(2n_GU)).
\]
Set $\varrho_y=[8n_G(1+m_yU)]^{-1}$. Then
$1/(2n_GU)-\varrho_y\geq3/(8n_GU)$. The guard threshold grows faster than
every power of $\log y$, so eventually
$3v^2/(8n_GU)\geq\log(2m_y)-2$. Hence
\begin{equation}\label{eq:readerB-correlated-reserve}
 \mathbb P\{\max_b\|X_b\|\geq v\}
       \leq e^2\exp(-\varrho_y v^2).
\end{equation}
The vector is the actual joint Gaussian. Correlation is allowed throughout;
only its marginal variance bound and the finite union bound are used.

Take the guard width
\[
 v=\frac{\operatorname{haze}(y,\beta)}2
   =\frac{\kappa_0}{2}\frac{\operatorname{phase}(y)}{g}.
\]
For the band $x/2\leq\beta\leq3x$, the scalar coupling bound gives
$v\geq\delta_0\,\operatorname{hot}(y)$ with a fixed $\delta_0>0$;
here $\operatorname{hot}(y)=\exp((\log\log y)^{3/2})$. By
\eqref{eq:readerB-room-dimension},
\[
 1+m_yU\leq (1+C_wU)(\log y)^{16}.
\]
For every fixed $p\geq0$, the ratio
\[
 \frac{\delta_0^2\exp(2(\log\log y)^{3/2})}
      {8n_G(1+C_wU)(\log y)^{16}(2+p)\log y}
\longrightarrow\infty.
\]
Indeed, writing $q=\log\log y$, the logarithm of the numerator growth
relative to the denominator is $2q^{3/2}-17q+O(1)$, which tends to
$+\infty$. Thus eventually
$\varrho_y v^2\geq(2+p)\log y$, and
\eqref{eq:readerB-correlated-reserve} gives
\begin{equation}\label{eq:readerB-tail-power}
 \mathbb P\{\max_b\|X_b\|\geq v\}
       \leq e^2 y^{-(2+p)}\leq x^{-p}.
\end{equation}

Let $\Gamma$ be the actual source veil, the product of smooth radial
cutoffs of the Lie-valued tape on each bond in the parasol. Each factor is
in $[0,1]$, is supported below the full vector guard threshold, and equals
one below half that threshold. Consequently
\[
 0\leq1-\Gamma(X)\leq
      \mathbf1_{\{\max_b\|X_b\|\geq v\}}.
\]
Integrating and using \eqref{eq:readerB-tail-power} proves
\[
 \mathbb E[1-\Gamma(X)]\leq x^{-p}.
\]
This is the lost-mask estimate for the same correlated Gaussian and its
actual finite room. It is not a product of single-coordinate probabilities.

\begin{lemma}[One-patch Gaussian guard]
\label{lem:readerB-mask}
For every fixed $p\geq0$, fixed patch width, and common palette entry bound
$U\geq1$, the mask loss on the actual joint Gaussian sample is at most
$x^{-p}$ for all sufficiently large $x$, uniformly in the period. In
particular, choose $p=1$ to obtain the mask loss $x^{-1}$ used in the
reader purchase.
\end{lemma}
\begin{proof}
The coordinate count is \eqref{eq:readerB-room-dimension}; the one-coordinate
Chernoff estimate, union bound and scale comparison establish
\eqref{eq:readerB-tail-power}. The pointwise cutoff comparison then proves
the expected mask loss.
\end{proof}

\begin{remark}[Boundary of this package]
The result proves the posed reader's nonlinear curvature/read estimates
and a power-small loss for its actual vector veil, conditional on the
same-pose Gaussian comparison, the paid local chart/ticket bounds, and the
fixed transition/window receipts. The source's sharper ShadePay input
uses its exact finite-dimensional Gaussian reserve bound. This unit does not
prove a different all-matched-palette socket assertion: that assertion
quantifies over a distinct family of matched palettes and requires its
separate local opacity estimate. Nor does this section produce the
selected history or the paid reader purchases.
\end{remark}
